\section{Related Work}
\label{sec:related}

\subsection{Graph Measures and Network Failures}

Freeman's account of centrality separated several ideas of what makes a node central, some
resting on a node's direct ties and others on the shortest paths that pass through it, and
gave them normalised forms that remain standard \cite{freeman1978}. Path-based measures became
practical on large graphs once Brandes published a fast algorithm for betweenness
\cite{brandes2001}. The graph library used in this study, NetworkX \val{version_networkx}
\cite{hagberg2008}, implements that algorithm.

Network science has studied how graphs respond to the loss of nodes. Albert, Jeong and
Barab\'asi examined the tolerance of complex networks to errors and to attacks
\cite{albert2000}, and Motter and Lai examined attacks on complex networks that work through
cascades \cite{motter2002}. In a running microservice system a failed service stays in the call
graph, and its effect appears as latency and errors in the services around it, which have to be
measured on the system itself.

Software engineering adopted graph measures early. Zimmermann and Nagappan computed network
measures on the dependency graph of a large software system and found that they helped predict
which components would contain defects \cite{zimmermann2008}, and Ma and colleagues built a
service dependency graph to analyse and test microservice systems \cite{ma2018}. The present
study shares their premise that a component's position in a dependency graph says something
about its risk. Its outcome is different: how far a run-time failure spreads under load, which
depends on the call pattern of the running system and not on the code alone.

\subsection{Fault Injection and Chaos Engineering}

Injecting faults into a running system is now an established way to learn how it behaves under
stress. Basiri and colleagues describe chaos engineering as practised at Netflix, where
controlled experiments on production systems build confidence that they tolerate failure
\cite{basiri2016}, and Basiri and colleagues later described automating chaos experiments in
production \cite{basiri2019}. For microservices, Gremlin offered systematic resilience testing by injecting
faults between services \cite{heorhiadi2016}, and Meiklejohn and colleagues brought fault
injection testing down to the level of individual services \cite{meiklejohn2021}.
Lineage-driven fault injection reasons backwards from how a correct result was produced to
choose the faults most likely to expose a bug \cite{alvaro2015}. Not every failure is a crash.
Huang and colleagues describe gray failure, in which a component keeps running while working
badly, as a central weakness of cloud-scale systems \cite{huang2017}; an injected network delay
is one controlled form of such a partial failure.

These tools search for bugs or confirm that fallbacks work. This study uses fault injection
for a narrower purpose. Two fault types are injected into each service in turn in a fixed,
balanced design, and the counts of degraded services become the outcome against which a
structural prediction is tested. The faults serve as a measuring instrument.

\subsection{Call Graphs from Traces}

Distributed tracing makes it possible to recover the service graph from what the system
actually does. Dapper set out the model most tracing systems still follow: a request carries a
trace context across service boundaries, each unit of work becomes a span with a parent, and a
request becomes a tree of spans \cite{sigelman2010}. A parent span in one service with a child
span in another gives a caller-to-callee edge, which is how the graphs in this study were
built. Graphs recovered this way depend on the traffic that produced them. Luo and colleagues
found the call graphs in Alibaba's production traces to be highly dynamic, changing from one
request to the next \cite{luo2021}. The benchmark graphs here are small and their workload is
fixed, but the same dependence on traffic reaches the predictor
(Section~\ref{sec:threat-scope}).

\subsection{Failure Diagnosis}

A large body of work starts from an anomaly that is already under way and tries to find its
source. Pinpoint addressed problem determination in large, dynamic Internet services
\cite{chen2002}. For microservices, Microscope
locates performance problems with causal graphs \cite{lin2018}, MicroRCA ranks candidate root
causes on a graph of the running system \cite{wu2020}, and Seer applies large-scale data
analysis to performance debugging in cloud microservices \cite{gan2019seer}. An industrial
survey documents how hard such diagnosis is in practice and contributes a benchmark with
injected faults \cite{zhou2021}, and Soldani and Brogi survey anomaly detection and
root-cause analysis for service-based cloud applications \cite{soldani2022}. These methods use
the graph together with live metrics after a problem has begun. The question here comes
earlier: whether graph structure alone, before any fault, indicates which services will do the
most damage when they fail. The direction of propagation reported in
Section~\ref{sec:res-direction} may interest diagnosis work, although no diagnosis method was
studied.

\subsection{Preregistration and Multiple Testing}

Preregistration separates the analyses planned before data are seen from those devised
afterwards, so that a reader can tell a test of a prediction from a search through the data
\cite{nosek2018}. With \val{conf_family_size} tests in each family, %[id:repl_family_size]
a correction for multiple testing was unavoidable. The false discovery rate was controlled
with the Benjamini-Hochberg procedure \cite{benjamini1995}, which bounds the expected share
of false positives among reported discoveries and gives up less power than a family-wise
correction.

\subsection{Research Gap}

Graph measures have been used to rank software components by defect risk
\cite{zimmermann2008}, network studies examine tolerance to errors, attacks and cascades
\cite{albert2000,motter2002}, fault-injection tools test whether a system tolerates
failure \cite{basiri2016,heorhiadi2016,meiklejohn2021}, and diagnosis methods use the graph
once an anomaly is under way \cite{wu2020,lin2018}. To the authors' knowledge, none of the
works cited here reports a test, fixed in advance, of whether a graph-derived ranking of
services anticipates failure propagation measured in a controlled campaign. This paper
reports such a test, including a result that came out null.
